\documentclass[a4paper]{standalone}
\usepackage{graphicx}
\usepackage{amsmath,amssymb}
\usepackage{hyperref}

\title{Efficient coding hypothesis}
\begin{document}
    \maketitle

\end{document}


\documentclass[a4paper]{article}

\usepackage{graphicx}
\usepackage{amsmath,amssymb}
\usepackage{minted}
\usepackage{multirow}
\usepackage{float}
\usepackage{todonotes}
\usepackage{forest}
\usepackage{xstring}
\usepackage{xcolor}
\usepackage[most]{tcolorbox}
\usepackage{xparse}
\usepackage{natbib}

\usetikzlibrary{graphs,graphdrawing}
\usegdlibrary{layered}

% for external graphviz
\input{/home/donkarlo/Dropbox/repo/nd_language_ai_project/src/nd_language_ai/formal/computer/programming/tex/latex/code_clips/external_graphviz.tex}

% for tags
\input{/home/donkarlo/Dropbox/repo/nd_language_ai_project/src/nd_language_ai/formal/computer/programming/tex/latex/parts/control_sequence/kind/semtag/semtag.tex}

% for links
\input{/home/donkarlo/Dropbox/repo/nd_language_ai_project/src/nd_language_ai/formal/computer/programming/tex/latex/code_clips/seehere.tex}

\title{}
\date{}
\author{Mohammad Rahmani}

\begin{document}
    \maketitle
    The efficient coding hypothesis was proposed by Horace Barlow in 1961 as a theoretical model of sensory neuroscience in the brain. Within the brain, neurons communicate with one another by sending electrical impulses referred to as action potentials or spikes. Barlow hypothesized that the spikes in the sensory system formed a neural code for efficiently representing sensory information. By efficient it is understood that the code minimized the number of spikes needed to transmit a given signal \seeherefooter{https://en.wikipedia.org/wiki/Efficient_coding_hypothesis}.
    
    Barlow proposes that sensory pathways actively transform incoming information rather than merely transmitting it. He describes sensory relays as mechanisms that can detect behaviorally important patterns or ``passwords'' in sensory input. Sensory relays may also dynamically control and filter information according to the current needs of the nervous system. The central proposal is that sensory systems reduce redundancy while preserving important information, producing a more efficient neural representation. This redundancy-reduction principle suggests that unusual or informative sensory events should produce stronger neural signals than common and predictable events \seeherefooter{barlow-1961-possible-principles-underlying-the-transformations-of-sensory-messages}.
    \bibliographystyle{plainnat}
    \bibliography{/home/donkarlo/Dropbox/repo/nd_shared_research/bibliography/bibliography.bib}
\end{document}